\documentclass{article}
\usepackage{amsmath, amssymb, amsthm}
\usepackage{hyperref}
\usepackage{geometry}
\geometry{margin=1in}

\title{Erd\H{o}s Problem \#201}
\date{Last updated: 2025-08-31}
\author{Source: erdosproblems.com}

\begin{document}
\maketitle

\noindent\textbf{Status:} OPEN \quad \textbf{No prize}

\noindent\textbf{Tags:} additive combinatorics, arithmetic progressions

\medskip
\noindent\textbf{Problem Statement:}

Let $G_k(N)$ be such that any set of $N$ integers contains a subset of size at least $G_k(N)$ which does not contain a $k$-term arithmetic progression. Determine the size of $G_k(N)$. How does it relate to $R_k(N)$, the size of the largest subset of $\{1,\ldots,N\}$ without a $k$-term arithmetic progression? Is it true that\[\lim_{N\to \infty}\frac{R_3(N)}{G_3(N)}=1?\]

\bigskip
\noindent\textbf{Remarks:}

First asked and investigated by Riddell \cite{Ri69}. It is trivial that $G_k(N)\leq R_k(N)$, and it is possible that $G_k(N) <R_k(N)$ (for example $G_3(5)=3$ and $R_3(5)=4$, and $G_3(14)\leq 7$ and $R_3(14)=8$).

Koml\'{o}s, Sulyok, and Szemer\'{e}di \cite{KSS75} have shown that $R_k(N) \ll_k G_k(N)$.

\bigskip
\noindent\textbf{References:}

[KSS75] Koml\'{o}s, J. and Sulyok, M. and Szemeredi, E., Linear problems in combinatorial number theory. Acta Math. Acad. Sci. Hungar. (1975), 113-121.
 
 [Ri69] Riddell, J., On sets of numbers containing no $l$ terms in arithmetic progression. Nieuw Arch. Wisk. (3) (1969), 204-209.

\bigskip
\noindent\small{Source: \url{https://www.erdosproblems.com/201}}
\end{document}
